\section{A functional consequence for the physical density}
\label{sec:hxc-first-variation}

In this section all coordinates and densities are physical, as in
\Cref{thm:main}.  Work on the real Banach space
$X=L^1(\R^3;\R)\cap L^3(\R^3;\R)$ with the sum norm and its potential
dual $\mathcal V$, paired by integration.  In addition to the kinetic
search \eqref{eq:ensemble-search}, define
\[
 F_{\mathrm{ens}}[m]
 =\inf_{\Gamma\mapsto m}\Tr((T+W)\Gamma),
 \qquad
 T=\sum_{i=1}^3-\frac12\Delta_i,
 \quad W=\sum_{i<j}\frac1{|r_i-r_j|}.
\]
Both searches run over all positive trace-one ensembles on
$\mathfrak H_3$, with form expectations, and take value $+\infty$ when
there is no finite-energy competitor.  Their common finite-energy
domain and ensemble Hartree--exchange--correlation functional are
\[
 \mathcal D=\Dom T_{\mathrm{ens}}=\Dom F_{\mathrm{ens}},
 \qquad
 E_{\mathrm{Hxc}}=F_{\mathrm{ens}}-T_{\mathrm{ens}}
 \quad\text{on }\mathcal D.
\]
Indeed, the Coulomb Hardy estimate from \Cref{sec:setup} gives
$0\le W\le\varepsilon T+C_\varepsilon$ as quadratic forms, also after
taking ensemble expectations.  Thus finite kinetic and interacting
internal energies are equivalent.  Such densities belong to $X$:
\eqref{eq:ensemble-density} and Cauchy--Schwarz give
$\int|\nabla\sqrt{n_\Gamma}|^2\le2T[\Gamma]$, and Sobolev together
with $\int n_\Gamma=3$ gives the assertion.  Mixtures and approximate
minimizers show that both searches are convex, so $\mathcal D$ is convex.
For either search $G$, write
\[
 \partial_{\mathcal V}G[n]
 =\left\{u\in\mathcal V:
 G[m]\ge G[n]+\int u(m-n)\ \text{for every }m\in X\right\}.
\]

\begin{corollary}[Absence of a full-domain local first variation]
\label{cor:hxc-first-variation}
For the physical ground-state density $n=n_\Psi$ of \Cref{thm:main},
$\partial_{\mathcal V}T_{\mathrm{ens}}[n]=\varnothing$.
Moreover, there is no $w\in\mathcal V$ such that, for every
$m\in\mathcal D$,
\begin{equation}
 \lim_{t\downarrow0}
 \frac{E_{\mathrm{Hxc}}[n+t(m-n)]-E_{\mathrm{Hxc}}[n]}{t}
 =\int_{\R^3}w(m-n).
 \label{eq:hxc-full-segment-variation}
\end{equation}
\end{corollary}

\begin{proof}
Proposition~\ref{prop:ensemble-certificate} supplies an attaining determinant for
the scaled kinetic search.  The inverse dilation in
Proposition~\ref{prop:scaling} gives a determinant $\Phi_{\rm phys}$ attaining
$T_{\mathrm{ens}}[n]$: it is a bijection of ensembles with the
corresponding densities and multiplies kinetic energies by $Z^2$.
If $u\in\partial_{\mathcal V}T_{\mathrm{ens}}[n]$, the subgradient
inequality makes $\Phi_{\rm phys}$ a minimizer of the form of $H_s(-u)$
over all normalized antisymmetric states in its form domain.
The form bound \eqref{eq:ensemble-standard-form-bound} applies to $u$,
so equality at the bottom of this closed form makes $\Phi_{\rm phys}$
a ground eigenstate, as in the proof of
Proposition~\ref{prop:potential-identification}.  Its density is $n$, contradicting
\Cref{thm:main}.

Let $v_{\rm ext}\in\mathcal V$ be the physical nuclear potential in
$H_C$.  The ground-state variational principle gives
$-v_{\rm ext}\in\partial_{\mathcal V}F_{\mathrm{ens}}[n]$;
in particular, $\Psi$ attains the interacting search at $n$.
Suppose now that $w$ satisfies \eqref{eq:hxc-full-segment-variation}.
Fix $m\in\mathcal D$, set $h=m-n$ and $n_t=n+th$.  Convexity of the
kinetic search and the interacting support inequality give, for
$0<t\le1$,
\begin{align*}
 T_{\mathrm{ens}}[m]-T_{\mathrm{ens}}[n]
 &\ge\frac{T_{\mathrm{ens}}[n_t]-T_{\mathrm{ens}}[n]}{t}\\
 &\ge-\int v_{\rm ext}h
       -\frac{E_{\mathrm{Hxc}}[n_t]-E_{\mathrm{Hxc}}[n]}{t}.
\end{align*}
Letting $t\downarrow0$ proves the subgradient inequality for
$-v_{\rm ext}-w$ and every $m\in\mathcal D$.  Outside $\mathcal D$ it
is automatic because $T_{\mathrm{ens}}=+\infty$.  Hence
$-v_{\rm ext}-w\in\partial_{\mathcal V}T_{\mathrm{ens}}[n]$, contrary
to the first assertion.
\end{proof}
